\section{Исходная программа}

\begin{center}
\lstinputlisting[language=C++,caption={Используемые include, c выбором под конкретную ОС},captionpos=b]{../src/os\_interaction/basic/defines.h}
\end{center}

\begin{center}
\lstinputlisting[language=C++,caption={Основной интерфейс взаимодействия с ОС, используемый в коде основной программы},captionpos=b]{../src/os\_interaction/basic/os.h}
\end{center}

\begin{center}
\lstinputlisting[language=C++,caption={Макросы для логирования},captionpos=b]{../src/os\_interaction/basic/logs.h}
\end{center}

\begin{center}
\lstinputlisting[language=C++,caption={Реализация взаимодействия с ОС POSIX},captionpos=b]{../src/os\_interaction/POSIX/os.cpp}
\end{center}

\begin{center}
\lstinputlisting[language=C++,caption={Реализация взаимодействия с ОС Windows},captionpos=b]{../src/os\_interaction/Windows/os.cpp}
\end{center}

\begin{center}
\lstinputlisting[language=C++,caption={Основная программа},captionpos=b]{../src/main.cpp}
\end{center}

\begin{center}
\lstinputlisting[language=C++,caption={Программа первого дочернего процесса},captionpos=b]{../src/child1.cpp}
\end{center}

\begin{center}
\lstinputlisting[language=C++,caption={Программа второго дочернего процесса},captionpos=b]{../src/child2.cpp}
\end{center}